\chapter{Wstęp}

Głębokie sieci neuronowe znacząco wpłynęły na rozwój dziedziny widzenia komputerowego (ang. \textit{Computer Vision}), osiągając wysoką skuteczność w zadaniach detekcji obiektów. Osiągnięcie takich rezultatów wymaga jednak dostępu do dużych, dokładnie oznaczonych zbiorów danych. W wielu rzeczywistych scenariuszach inżynieryjnych i przemysłowych pozyskanie takich zbiorów jest niemożliwe, co stwarza istotną barierę we wdrażaniu algorytmów sztucznej inteligencji. Niniejsza praca podejmuje ten problem, analizując metody poprawy zdolności generalizacji modeli wizyjnych w warunkach ograniczonej dostępności danych treningowych.

\section{Tło problemu}

Małe zbiory danych stanowią istotne wyzwanie w procesie trenowania współczesnych modeli detekcji obiektów. Architektury takie jak modele z rodziny YOLO czy oparte na transformerach sieci RT-DETR posiadają dużą liczbę parametrów (od kilku do kilkudziesięciu milionów). W przypadku trenowania ich na tzw. małym zbiorze danych (ang. \textit{low-data regime}) sieci te są podatne na przeuczenie (ang. \textit{overfitting}). Zamiast uczyć się ogólnych, niezmiennych cech obiektów docelowych, model zaczyna zapamiętywać specyficzne artefakty, szum obrazu oraz tekstury tła. W efekcie znacząco obniża się jego zdolność do generalizacji na nowe, niewidziane wcześniej dane, co może prowadzić do występowania fałszywie pozytywnych detekcji w miejscach, gdzie znajduje się wyłącznie tło.

Praktyczną inspiracją do podjęcia tego tematu był projekt realizowany we współpracy z Wydziałem Geoinżynierii, Górnictwa i Geologii Politechniki Wrocławskiej. Jego celem było opracowanie prototypu systemu do zautomatyzowanej detekcji i zliczania worków z piaskiem na podstawie zdjęć lotniczych pozyskiwanych z bezzałogowych statków powietrznych (UAV), potocznie określanych jako drony. Informacje o liczbie oraz rozmieszczeniu worków mogą wspierać ocenę stopnia zabezpieczenia wałów przeciwpowodziowych, identyfikację potencjalnych miejsc osłabienia zabezpieczeń oraz usprawniać planowanie działań logistycznych służb odpowiedzialnych za ochronę przeciwpowodziową.

Zadanie to okazało się wymagające. Po pierwsze, obiekty tego typu występują niemal wyłącznie w rzadkich sytuacjach kryzysowych, takich jak zagrożenie powodziowe, co znacząco ogranicza możliwość pozyskiwania nowych danych i rozbudowy zbioru treningowego. Po drugie, charakterystyka wizualna analizowanych scen — obejmująca gęste upakowanie obiektów, ich wzajemne zasłanianie oraz złożone tło o cechach zbliżonych do wyglądu worków — utrudnia proces uczenia modeli detekcyjnych. Ze względu na niewielką objętość pierwotnego zbioru danych pojawiła się potrzeba zbadania i porównania zaawansowanych metod wzbogacania danych treningowych, które mogłyby ograniczyć zjawisko przeuczenia i poprawić zdolność generalizacji modeli.


\section{Cel i zakres pracy}

Głównym celem badawczym niniejszej pracy jest analiza wpływu metod wzbogacania danych treningowych na zdolność generalizacji modeli detekcji obiektów trenowanych w warunkach ograniczonej dostępności danych. W ramach badań porównano dwa podejścia należące do nurtu data-centric AI: iteracyjne uczenie półnadzorowane (ang. \textit{Semi-Supervised Learning}, SSL) wykorzystujące pseudoetykietowanie oraz syntetyczną augmentację danych polegającą na generowaniu próbek negatywnych przy użyciu modeli dyfuzyjnych. Skuteczność obu metod zweryfikowano eksperymentalnie dla trzech architektur detekcyjnych: YOLOv8n, YOLOv26n oraz RT-DETR-L.
\todo[inline]{Wspólnie poprawić cele, koniecznie odwołać się do nich również w analizie wyników eksperymentalnych, a nie wyłączenie posumowaniu...}

Dodatkowym wyzwaniem implementacyjnym podjętym w ramach pracy było zaprojektowanie i wdrożenie zautomatyzowanego środowiska eksperymentalnego umożliwiającego powtarzalne trenowanie, ewaluację oraz porównywanie modeli. Środowisko zostało oparte na technologiach konteneryzacji i przystosowane do pracy na infrastrukturze obliczeniowej NVIDIA DGX A100, co umożliwiło realizację wieloetapowych eksperymentów obejmujących pętle uczenia Nauczyciel–Uczeń.

\subsection{Hipoteza badawcza}

Hipoteza badawcza zakłada, że wykorzystanie dodatkowych danych, zarówno rzeczywistych jak i syntetycznych, może poprawić zdolność generalizacji modeli detekcji obiektów trenowanych na ograniczonym zbiorze danych. Przyjęto również, że skuteczność poszczególnych metod może różnić się w zależności od zastosowanej architektury modelu.

\section{Wykorzystanie narzędzi sztucznej inteligencji}

W procesie przygotowywania niniejszej pracy dyplomowej skorzystano ze wsparcia modeli językowych, w szczególności modelu Google Gemini 3.1 Pro Preview poprzez platformę Google AI Studio. Narzędzia te zostały użyte jako pomoc w pracach redakcyjnych, korekcie stylistycznej i gramatycznej tekstu, a także przy formatowaniu środowiska składu dokumentów \LaTeX. Ponadto wsparcie sztucznej inteligencji wykorzystano podczas etapu programowania zautomatyzowanego środowiska eksperymentalnego, m.in. do optymalizacji oraz debugowania skryptów w języku Python.

\section{Struktura pracy}

Oprócz niniejszego rozdziału wprowadzającego praca składa się z pięciu głównych części. Rozdział 2 przedstawia podstawy teoretyczne głębokich sieci neuronowych wykorzystywanych w zadaniach widzenia komputerowego. Omówiono w nim architektury detekcji obiektów z rodzin YOLO oraz RT-DETR, a także podstawowe metryki stosowane do oceny jakości modeli.

Rozdział 3 zawiera przegląd literatury poświęconej problemowi uczenia modeli w warunkach ograniczonej dostępności danych. Przedstawiono w nim wyzwania charakterystyczne dla małych zbiorów treningowych oraz omówiono współczesne podejścia wykorzystujące uczenie półnadzorowane i modele generatywne do poprawy jakości detekcji.

W rozdziale 4 zaprezentowano metodykę badań oraz szczegóły implementacyjne. Opisano wykorzystane środowisko eksperymentalne, sposób przygotowania i podziału danych, procedury pseudoetykietowania oraz proces generowania syntetycznych obrazów wykorzystywanych podczas augmentacji.

Rozdział 5 przedstawia wyniki przeprowadzonych eksperymentów wraz z ich analizą. Porównano skuteczność badanych metod dla poszczególnych architektur, analizując wpływ zastosowanych technik na proces uczenia, poziom przeuczenia oraz jakość predykcji.

Pracę zamyka rozdział 6 zawierający podsumowanie uzyskanych rezultatów, ocenę postawionych hipotez badawczych oraz propozycje dalszych kierunków rozwoju przedstawionych rozwiązań. Po części zasadniczej zamieszczono bibliografię oraz dodatek zawierający rozszerzony zestaw wizualizacji wyników uzyskanych podczas eksperymentów.

